\documentclass[10pt]{article}
\providecommand{\FIGDIR}{fig}
\newcommand{\DIGESTDATE}{6 September 2026}
% ===================== PM Research Watch - shared preamble =====================
\usepackage[T1]{fontenc}
\usepackage[utf8]{inputenc}
\usepackage{amsmath,amssymb}
\usepackage{tgpagella}
\usepackage{tgheros}
\renewcommand{\sfdefault}{qhv}
\renewcommand{\ttdefault}{lmtt}
\usepackage{microtype}
\usepackage[english]{babel}

\usepackage[a4paper,top=17mm,bottom=17mm,left=17mm,right=17mm,
            headsep=6mm,footskip=10mm]{geometry}
\usepackage{graphicx}
\usepackage{xcolor}
\usepackage{colortbl}
\usepackage{tikz}
\usetikzlibrary{calc}
\usepackage[most]{tcolorbox}
\usepackage{titlesec}
\usepackage{fancyhdr}
\usepackage{booktabs}
\usepackage{array}
\usepackage{longtable}
\usepackage{enumitem}
\usepackage{multicol}
\usepackage{ragged2e}
\usepackage[hidelinks]{hyperref}

% ------------------------------------------------------------------ palette
\definecolor{Ink}    {HTML}{14262E}
\definecolor{Deep}   {HTML}{0F4C5C}
\definecolor{Teal}   {HTML}{1F7A8C}
\definecolor{Sky}    {HTML}{4FA3C4}
\definecolor{Sage}   {HTML}{6FA287}
\definecolor{Amber}  {HTML}{E3A008}
\definecolor{Clay}   {HTML}{B5643C}
\definecolor{Coral}  {HTML}{D1495B}
\definecolor{Violet} {HTML}{6A5B9E}
\definecolor{Slate}  {HTML}{4A6572}
\definecolor{Paper}  {HTML}{FBF9F5}
\definecolor{Mist}   {HTML}{EEEAE2}
\definecolor{Rule}   {HTML}{D9D3C9}

\pagecolor{Paper}
\color{Ink}

\hypersetup{colorlinks=true, linkcolor=Deep, urlcolor=Teal, citecolor=Teal,
            pdftitle={PM Research Watch}, pdfauthor={Automated daily digest}}

% ------------------------------------------------------------------ headings
\titleformat{\section}
  {\sffamily\bfseries\large\color{Deep}}{}{0pt}
  {}[{\vspace{1.2mm}\color{Amber}\titlerule[1.1pt]}]
\titlespacing*{\section}{0pt}{5.5mm}{2.8mm}

\titleformat{\subsection}
  {\sffamily\bfseries\normalsize\color{Teal}}{}{0pt}{}
\titlespacing*{\subsection}{0pt}{3.6mm}{1.4mm}

\setlist[itemize]{leftmargin=4.2mm, itemsep=1.1mm, topsep=1.3mm, parsep=0pt,
                  label=\textcolor{Teal}{\footnotesize$\blacktriangleright$}}

% ------------------------------------------------------------------ header/footer
\pagestyle{fancy}
\fancyhf{}
\renewcommand{\headrulewidth}{0pt}
\renewcommand{\footrulewidth}{0pt}
\fancyhead[L]{\sffamily\scriptsize\color{Slate}PM RESEARCH WATCH\;\textcolor{Rule}{\textbar}\;\DIGESTDATE}
\fancyhead[R]{\sffamily\scriptsize\color{Slate}Particulate matter monitoring \& health effects}
\fancyfoot[C]{\sffamily\scriptsize\color{Slate}\thepage}

% ------------------------------------------------------------------ components
% metric tile
\newcommand{\metric}[3]{%
  \begin{tcolorbox}[enhanced, colback=white, colframe=Rule, boxrule=0.5pt,
    arc=1.6mm, left=2.2mm, right=2.2mm, top=1.8mm, bottom=1.8mm,
    borderline west={1.7mm}{0pt}{#3}, nobeforeafter, valign=center,
    height=17mm, width=\linewidth]
    \RaggedRight
    {\sffamily\bfseries\fontsize{17}{17}\selectfont\color{#3}#1}\\[0.5mm]
    {\sffamily\fontsize{6.5}{7.4}\selectfont\color{Slate}#2}
  \end{tcolorbox}}

% pull-quote / key insight box
\newtcolorbox{keybox}[1][]{enhanced, colback=Mist, colframe=Mist,
  arc=1.6mm, left=3.2mm, right=3.2mm, top=2.4mm, bottom=2.4mm,
  borderline west={1.7mm}{0pt}{Amber}, #1}

% subtopic band
\newcommand{\band}[2]{%
  \vspace{3.2mm}\noindent
  \begin{tcolorbox}[enhanced, colback=#2, colframe=#2, arc=1.2mm,
    left=2.6mm, right=2.6mm, top=1.1mm, bottom=1.1mm, boxrule=0pt,
    width=\linewidth, nobeforeafter]
    {\sffamily\bfseries\small\color{white}#1}
  \end{tcolorbox}\vspace{1.6mm}}

% one paper entry: {short cite}{title}{journal}{doi}{body}
\newcommand{\paperentry}[5]{%
  \par\vspace{1.9mm}\noindent
  {\sffamily\bfseries\footnotesize\color{Ink}#1}\;%
  {\sffamily\fontsize{7}{8}\selectfont\color{Slate}#3}\\[0.4mm]
  {\footnotesize\itshape\color{Deep}#2}\\[0.6mm]
  {\footnotesize #4}\\[0.3mm]
  {\fontsize{6.6}{7.6}\selectfont\color{Slate}\texttt{doi:}\,\href{https://doi.org/#5}{\texttt{\detokenize{#5}}}}
  \par}

% caption style
\newcommand{\figcap}[1]{\par\vspace{1.0mm}
  {\fontsize{7.2}{8.6}\selectfont\color{Slate}\textsf{\textbf{Fig.}}\; #1\par}\vspace{1.5mm}}

\providecommand{\permil}{\textperthousand}

\setlength{\parindent}{0pt}
\setlength{\parskip}{0pt}
\renewcommand{\arraystretch}{1.18}


\begin{document}

% ==================================================================== MASTHEAD
\thispagestyle{empty}
\begin{tikzpicture}[remember picture, overlay]
  \fill[Deep] (current page.north west) rectangle
      ($(current page.north east)+(0,-67mm)$);
  \fill[Amber] ($(current page.north west)+(0,-67mm)$) rectangle
      ($(current page.north east)+(0,-68.4mm)$);
  \foreach \x/\y/\r in {22/12/0.9, 41/26/1.5, 63/9/0.7, 88/31/1.1, 112/17/1.9,
                        138/29/0.8, 159/13/1.3, 178/54/1.0, 196/21/1.6,
                        30/44/0.6, 74/57/1.2, 125/48/0.9, 168/6/0.8, 191/59/0.7,
                        52/18/1.0, 100/60/0.6, 148/20/1.1, 66/49/1.4, 205/45/0.8}
    \fill[Sky, opacity=0.30] ($(current page.north west)+(\x mm,-\y mm)$) circle (\r mm);
\end{tikzpicture}

\vspace*{4mm}
{\sffamily\fontsize{7.6}{9}\selectfont\color{Sky}%
 AUTOMATED DAILY LITERATURE DIGEST\;\textbullet\;PhD RESEARCH WATCH\par}
\vspace{2.2mm}
{\sffamily\bfseries\fontsize{27}{29}\selectfont\color{white}%
 Particulate Matter\par}
{\sffamily\fontsize{27}{29}\selectfont\color{Sky}%
 Monitoring \& Health Effects\par}
\vspace{5.0mm}
{\sffamily\fontsize{8.4}{10}\selectfont\color{white}%
 Issue 2026-09-06\;\;\textcolor{Amber}{\textbar}\;\;Entry window 6 September 2026\;\;%
 \textcolor{Amber}{\textbar}\;\;15 records screened in scope\par}
\vspace{18mm}

% -------------------------------------------------------------- metric strip
\noindent
\begin{minipage}[t]{0.190\linewidth}\metric{15}{RECORDS IN SCOPE\\(21 EXCLUDED)}{Deep}\end{minipage}\hfill
\begin{minipage}[t]{0.190\linewidth}\metric{10}{EFFECT ESTIMATES\\EXTRACTED}{Teal}\end{minipage}\hfill
\begin{minipage}[t]{0.190\linewidth}\metric{10}{WITH A HUMAN\\HEALTH ENDPOINT}{Sky}\end{minipage}\hfill
\begin{minipage}[t]{0.190\linewidth}\metric{4}{ON THE SENSING /\\EXPOSURE AXIS}{Amber}\end{minipage}\hfill
\begin{minipage}[t]{0.190\linewidth}\metric{0}{TRIAL UPDATES\\(SUNDAY ZERO)}{Clay}\end{minipage}

\vspace{6mm}

% -------------------------------------------------------------- signal of day
\section*{\sffamily\bfseries\color{Deep}Signal of the day}
\vspace{-1.5mm}{\color{Amber}\rule{\linewidth}{1.1pt}}\vspace{2.5mm}

\begin{keybox}
{\sffamily\bfseries\small\color{Deep}The PM$_{2.5}$ concentration--response slope is steepening
as concentrations fall --- and it has already clawed back a fifth of China's clean-air
mortality benefit.}\\[1.4mm]
\footnotesize
A multi-county time-series across the Sichuan Basin (\textit{Int Arch Occup Environ Health},
2014--2022, extended two-stage design with longitudinal meta-regression) finds the excess
risk of cardiorespiratory mortality per 10\,$\mu$g/m$^3$ PM$_{2.5}$ moved
\textbf{non-monotonically} across the phases of the Clean Air Action:
\textbf{0.7\% (95\% CI 0.6--0.9) in 2014--2016, 0.4\% (0.3--0.6) in 2017--2019, then
1.4\% (1.2--1.6) in 2020--2022} ($P<0.001$ for the temporal pattern). Decomposed against a
counterfactual, falling concentrations avoided \textbf{52{,}831} cardiorespiratory deaths
while the drift in the exposure--response function \emph{added} \textbf{10{,}041} --- so
roughly 19\% of the achieved benefit was given back, concentrated in those aged 65 and over.
The mechanism is unresolved and the paper does not claim one; changing composition at lower
mass, a shifting susceptible fraction, and residual confounding by the 2020--2022 period
itself are all live. The instrumentation consequence is sharp and independent of which
explanation wins: as the mean falls, the health-relevant contrast shrinks while the
per-unit risk grows, so a network's useful resolution has to improve just to hold
epidemiological power constant. A monitor calibrated to be adequate at 60\,$\mu$g/m$^3$ is
not adequate at 25, and mass-only calibration cannot detect a composition-driven slope
change at all.
\end{keybox}

\vspace{3mm}

% -------------------------------------------------------------- exec brief
\section*{\sffamily\bfseries\color{Deep}What changed today}
\vspace{-1.5mm}{\color{Amber}\rule{\linewidth}{1.1pt}}\vspace{2.5mm}

\footnotesize
\begin{itemize}
\item \textbf{PubMed deposits on Sundays --- the 5 September zero was not a weekend rule.}
This window is a Sunday and PubMed returned \textbf{9 health-axis and 1 sensing-axis} hits
to the harvester, and 10 and 13 to the connector. Saturday's both-ways-verified zero was
therefore a property of that day, not of weekends. \textbf{ClinicalTrials.gov is the
opposite case} and its weekend zero \emph{is} structural: the no-term control returns
\textbf{0 studies updated registry-wide} on 6 September, as on 5 September and on
29/30 August.

\item \textbf{Fifteen records and ten effect estimates --- the forest plot returns after a
one-day gap.} 5 September carried no estimate with a confidence interval at all. Today
eight of the ten come from just two papers, so the panel shows within-study precision far
more than between-study agreement, and is captioned as such.

\item \textbf{Solid fuel outranks ambient PM inside the mixture.} In CHARLS
($n=10{,}273$ for cardiometabolic multimorbidity), quantile g-computation puts every
co-exposure up one quantile together for an \textbf{OR of 1.650 (1.375--1.979)}, and
identifies \textbf{indoor solid-fuel use as the primary contributor} --- ahead of PM$_1$,
PM$_{2.5}$, PM$_{10}$ and heatwave. Household combustion keeps outperforming outdoor mass
in mixture weights whenever both are measured together.

\item \textbf{Day-to-day \emph{change} in PM$_{2.5}$, not level, predicts infarction.}
The Zhongshan case-crossover (1{,}795 admissions) reports \textbf{+7.89\%
(0.86--15.40) odds per 10\,$\mu$g/m$^3$ of lag0--1 increment}, rising to 17.90\% for
ST-elevation infarction. A first-difference exposure metric is cheap for a sensor network
to compute and is insensitive to the constant offset that dominates low-cost PM error ---
the rare case where an uncalibrated node still carries the signal.

\item \textbf{A low-cost sensor is the exposure instrument in a health study, and the
result is a warning.} The Douala survey (1{,}721 individuals, 344 households) drives an
instrumental Logit from an OC~300 Dust Particle Laser and reports effects of
\textbf{0.231 (PM$_{10}$), 0.122 (PM$_5$) and 0.018 (PM$_{2.5}$)} --- monotonically
\emph{decreasing} with decreasing size, which inverts the usual deposition argument. An
uncorrected optical instrument apportioning mass across bins it cannot resolve is the
parsimonious explanation, and it is what a co-location arm would have settled.

\item \textbf{A clean null: exercise benefit is not modified by industrial emissions.}
CLHLS, 16{,}954 adults aged 65+, six years: habitual exercise \textbf{HR 0.835
(0.773--0.902)}, county industrial SO$_2$ \textbf{HR 1.001 (0.874--1.146)}, and
\textbf{no significant interaction} on any of four pollution indicators. County-level
soot and SO$_2$ tonnage is a weak exposure surrogate, so this bounds the design more than
the hypothesis.

\item \textbf{The connector's sensing axis is now mostly noise.} 13 hits, \textbf{10
rejected outright} and only two new to this issue. Free-text \texttt{"remote sensing"} and
\texttt{"exposure assessment"} pulled in phycotoxin hydrography and a football-heading
index; particle terms must be ANDed in.
\end{itemize}

% ==================================================================== ANALYTICS
\section*{\sffamily\bfseries\color{Deep}Corpus analytics}
\vspace{-1.5mm}{\color{Amber}\rule{\linewidth}{1.1pt}}\vspace{3mm}

\noindent\includegraphics[width=\linewidth]{\FIGDIR/f1_subtopics.png}
\figcap{Subtopic assignment of the 15 in-scope records. Eight of the ten clusters are
occupied; \textit{Reproductive \& developmental} and \textit{Other clinical endpoints} are
empty. \textit{Cardiovascular \& metabolic} leads with three, and no cluster exceeds three
--- an unusually flat distribution for an issue of this size, and the reason no single
theme carries the day.}

\vspace{2mm}
\noindent\includegraphics[width=\linewidth]{\FIGDIR/f2_design_endpoint.png}
\figcap{Left --- study architecture. Observational cohorts account for 4 of 15 and
measurement campaigns for 3; the two \textit{Modelling / inventory} entries are the SAQIN
inference network and the California wine-grape regression. Right --- health endpoint.
\textbf{Ten records carry a human health endpoint}, the inverse of 5 September, when only
two did. Respiratory leads at three; the five without a health endpoint are the two
sensing records, the exposome review, the crop-yield preprint and the in-vitro
toxicology.}

\vspace{2mm}
\noindent\includegraphics[width=\linewidth]{\FIGDIR/f5_forest.png}
\figcap{All ten quantitative estimates in the corpus, on a common logarithmic axis.
\textbf{The exposure contrasts are not mutually comparable} --- four are quantile
g-computation mixture ORs in which every co-exposure moves together, two are per-10\,%
$\mu$g/m$^3$ day-to-day PM$_{2.5}$ increments, two are binary solid-fuel contrasts, and two
are non-particulate (habitual exercise, county SO$_2$ tonnage). Read direction and
precision only. Eight of the ten come from two papers, so the visual clustering is
within-study, not corroboration. The ST-elevation subtype estimate is plotted as a point
because the source reports no interval for it.}

\vspace{2mm}
\noindent\includegraphics[width=\linewidth]{\FIGDIR/f3_metric_geography.png}
\figcap{Left --- particle metric. \textit{PM (unspeciated) / emissions} is the largest bin
at 7, which overstates the field's precision: it holds county SO$_2$ and soot tonnage,
occupational coal dust by job history, dosed reference materials and self-reported cooking
fuel. Only 8 of 15 records measure a size-resolved particle metric at all. Right ---
geographic provenance. China supplies 5. The \textit{Global / multi-region} bar is a
fallback bucket, not a finding: one genuine multi-region review plus two records whose
setting is unstated and one bench study.}

\vspace{2mm}
\noindent\includegraphics[width=\linewidth]{\FIGDIR/f4_heatmap.png}
\figcap{Subtopic $\times$ study-architecture cross-tabulation. Twelve of 56 cells are
occupied and none holds more than two records, so this is a maximally dispersed day.
The structural point is that every two-record subtopic \emph{except} neuro spreads its
records across two different architectures. Note also that \textit{Exposure assessment}
is not siloed today: the Douala record sits in that row and carries a respiratory
endpoint, which is the linkage this cross-tabulation usually shows missing.}

\vspace{2mm}
\noindent\includegraphics[width=\linewidth]{\FIGDIR/f6_lifecourse.png}
\figcap{Life-course windows addressed; counts are non-exclusive. Three windows carry two
records each; the in-utero and adolescent windows are empty, and the reproductive and
developmental cluster is absent from this issue entirely.}

% ==================================================================== DIGEST
\section*{\sffamily\bfseries\color{Deep}Paper-level digest}
\vspace{-1.5mm}{\color{Amber}\rule{\linewidth}{1.1pt}}\vspace{1mm}

{\fontsize{7.4}{9}\selectfont\color{Slate}Ordered by cluster. Each entry gives the
methodological core and the finding that matters, not an abstract paraphrase.\par}

\band{Cardiovascular \& metabolic\hfill 3 records}{Teal}

\paperentry{Sichuan Basin CAA 2026}
{Temporal changes in short-term PM$_{2.5}$--cardiorespiratory mortality associations under
the Clean Air Action}
{Int Arch Occup Environ Health}
{Multi-county time-series, Sichuan Basin 2014--2022, extended two-stage design with
county-specific estimates pooled and longitudinal meta-regression on annual excess risk.
ER per 10\,$\mu$g/m$^3$ traced \textbf{0.7\% (0.6--0.9) $\to$ 0.4\% (0.3--0.6) $\to$
1.4\% (1.2--1.6)} across the three policy periods, $P<0.001$ for the pattern; scenario
decomposition gives 52{,}831 deaths averted by concentration decline against 10{,}041 added
by the risk change, and the late rise held in both sexes and was larger at ages 65+.
The threat to the inference is that the final period is 2020--2022: pandemic-era changes in
care-seeking, mortality coding and time--activity are not separable from a true slope change
in this design. Still the most consequential result here --- a falling mean does not imply a
falling per-unit risk, and networks must resolve smaller contrasts as policy succeeds.}
{10.1007/s00420-026-02237-z}

\paperentry{Zhongshan AMI 2026}
{Daily \emph{changes} in PM$_{2.5}$ and acute myocardial infarction hospitalisation}
{Front Public Health}
{Individual-level time-stratified case-crossover, 1{,}795 AMI admissions to one tertiary
hospital in Zhongshan 2020--2024, 5{,}877 control days matched on weekday within calendar
month; residential PM$_{2.5}$ from the CHAP 1\,km product; conditional logistic regression
adjusted for temperature, humidity and holidays. A 10\,$\mu$g/m$^3$ lag0--1 \emph{increment}
gave \textbf{+7.89\% (0.86--15.40)} odds, with 22.58\% in women, 22.16\% at ages $>$65,
10.85\% in the cold season and 17.90\% for STEMI; age modified the effect ($P_{int}<0.05$).
Robust to co-pollutants except NO$_2$, which is the honest weakness --- traffic co-exposure
is not excluded. Single-hospital catchment limits generalisability, and the 1\,km modelled
field is assigned to address, not measured. The transferable idea for sensing is the
first-difference metric: it cancels the additive bias that dominates low-cost PM error.}
{10.3389/fpubh.2026.1911027}

\paperentry{CHARLS coexposure 2026}
{Indoor solid fuels, outdoor PM and heatwave vs.\ cardiometabolic multimorbidity}
{Int Arch Occup Environ Health}
{Five CHARLS waves 2011--2020, four disease cohorts (CMM $n=10{,}273$; diabetes 9{,}893;
heart disease 9{,}536; stroke 10{,}466), time-varying Cox plus a cumulative risk index and
quantile g-computation. Single-exposure: solid fuel \textbf{HR 1.118 (1.016--1.231)} for
heart disease and \textbf{1.149 (1.007--1.310)} for stroke; PM$_1$, PM$_{2.5}$ and PM$_{10}$
were the most consistent risk factors ($P<0.001$). Mixture ORs: \textbf{CMM 1.650
(1.375--1.979), diabetes 1.675 (1.475--1.903)}, with solid fuel carrying the largest weight.
Fuel use is self-reported and correlates with income, education and rurality, so the weight
assigned to it is the estimate most exposed to residual confounding. Outdoor PM is modelled,
not measured, at residence.}
{10.1007/s00420-026-02232-4}

\band{Neuro / mental health\hfill 2 records}{Deep}

\paperentry{ENVIRONAGE eye-track 2026}
{Early-life air pollution and oculomotor markers of cognitive workload}
{EBioMedicine}
{196 children aged 9--11 from the ENVIRONAGE birth cohort completed eye-tracking paradigms
2021--2023; residential BC, NO$_2$ and PM$_{2.5}$ from a high-resolution spatiotemporal
model combining fixed-site monitors, land use and dispersion; multiple linear regression
then BKMR distributed-lag models across windows. After multiple-testing correction the
surviving associations are \emph{recent} exposure with saccadic outcomes --- per IQR of
last-month BC \textbf{--10.97$^\circ$/s (--18.67, --3.26)} and NO$_2$ \textbf{--10.81
(--18.13, --3.49)} saccade velocity --- plus first-trimester and first-1000-days windows on
the same outcome. Pupil diameter, fixation count and attention-bias outcomes did not
survive. $n=196$ against many windows and pollutants means the surviving hits are the ones
to trust and the null ones are underpowered rather than negative. Notable for sensing: the
effect attaches to last-month exposure, a window a personal or neighbourhood monitor can
actually measure.}
{10.1016/j.ebiom.2026.106469}

\paperentry{ABCD amygdala 2026 (preprint)}
{Air pollution, early-life adversity and amygdala volume in preadolescents}
{bioRxiv (preprint)}
{3{,}601 ABCD Study children aged 9--10, annual residential PM$_{2.5}$ total mass plus 15
components, linear mixed-effects adjusted for sociodemographics, co-pollutants and imaging
covariates. Adversity predicted general, internalising and externalising symptoms;
\textbf{PM$_{2.5}$ moderated the adversity--externalising association}, with the highest
symptoms in children high on both, and smaller basolateral paralaminar volumes tracked
symptoms within the high-adversity group. Cross-sectional, so the mediation reading the
framing invites is not supported; annual average exposure at residence cannot speak to the
timing the developmental argument needs; and with 16 exposure terms the component results
want a correction the abstract does not describe. Not peer reviewed. The component panel is
the part worth watching --- it is the kind of speciated contrast that mass-only networks
cannot support.}
{10.64898/2026.08.31.748438}

\band{Respiratory \& allergic\hfill 2 records}{Sage}

\paperentry{COPD anaemia metals 2026}
{Air-pollution-associated intracellular metals, immune gene expression and erythrocyte
indices}
{Int J Chron Obstruct Pulmon Dis}
{61 COPD patients and 10 controls; annual PM$_{10}$, PM$_{2.5}$, NO$_2$ and NO$_x$ by land
use regression; metals by ICP-MS and single-cell ICP-MS in PBMCs and red cells; PBMC
RNA-seq. PM$_{2.5}$ was associated with higher RBC count ($\beta=0.0314\times10^6/\mu$L,
0.0066--0.0561), red cell distribution width, and with intracellular chromium, zinc and
cadmium; those metals in turn tracked \textit{ALPL}, \textit{CXCR1} and \textit{PI3}
expression, and anaemic patients showed a distinct signature. Ten controls against 61 cases
and a cross-sectional design make every link correlational, and LUR annual means cannot
support a claim about cellular metal kinetics. Read it as an assay-development result:
single-cell ICP-MS as an internal-dose readout is the transferable piece, and it is the
kind of endpoint a speciating monitor could eventually be validated against.}
{10.2147/COPD.S626348}

\paperentry{Bengaluru COPD 2026}
{Prevalence and risk factors of COPD in an urban poor locality}
{Indian J Community Med}
{Community-based cross-sectional survey, adults over 30 in an urban poor locality of
Bengaluru, PPS sampling to a calculated 1{,}560, spirometry-confirmed diagnosis.
Prevalence \textbf{10.38\%} overall, 14.4\% in men and 7.1\% in women, with 74.06\% at
moderate-to-severe stage; multivariate logistic regression retained age, smoked and
smokeless tobacco, fuel type, biomass exposure, environmental tobacco smoke and indoor air
pollution. No air measurement of any kind --- every exposure is a questionnaire item, so
the ``indoor air pollution'' association is a self-report contrast and the effect sizes are
not reported in the abstract. Value here is the spirometry-confirmed prevalence in a
population that ambient monitoring networks systematically under-serve.}
{10.4103/ijcm.ijcm_145_25}

\band{Exposure assessment \& modelling\hfill 2 records}{Amber}

\paperentry{Douala households 2026}
{Particulate concentration and household respiratory health in urban Cameroon}
{Spat Spatiotemporal Epidemiol}
{On-site campaign in Douala with an OC~300 Dust Particle Laser --- a low-cost portable
optical sensor --- paired with a survey of 1{,}721 individuals in 344 households sited near
the measurement points; instrumental Logit to address endogeneity of exposure. Reported
effects on respiratory disease were \textbf{0.231 for PM$_{10}$, 0.122 for PM$_5$ and 0.018
for PM$_{2.5}$}, and unsealed roads and traffic density dominated the PM levels. The size
ordering is the problem: effect falls monotonically as particles get smaller, the reverse
of what deposition physics predicts, and the most economical explanation is an uncorrected
nephelometric instrument distributing mass across bins it cannot separate --- no
co-location against a reference monitor is described, and no humidity correction. Included
because it is a rare health study built on a low-cost node, and because it shows precisely
what the missing co-location arm costs.}
{10.1016/j.sste.2026.100836}

\paperentry{Paediatric exposome 2026}
{Assessing environmental exposures in children: toward a comprehensive exposome approach}
{Pediatr Res}
{Narrative review of exposure-assessment tooling for paediatrics --- questionnaires,
biomarkers, wearable and fixed sensors, geospatial indicators and pooled databases ---
alongside longitudinal cohort design, multi-omics integration and methods for mixtures and
critical windows. No systematic search protocol, no pooled estimate, and the closing
argument is professional advocacy rather than evidence. Useful only as a citation anchor
for the claim that single-pollutant assessment is the limiting factor in child
environmental epidemiology, and for its framing of sensors as one instrument among several
rather than the exposure itself.}
{10.1038/s41390-026-05460-z}

\band{Sensing, forecasting \& instrumentation\hfill 2 records}{Amber}

\paperentry{SAQIN 2026}
{Long-short-term decoupled residual learning for high-resolution spatio-temporal air
quality inference}
{Neural Netw}
{Inference --- not forecasting --- of high-resolution concentration fields from sparse
reference stations. The model splits a static branch (semantic segmentation, remote-sensing
context, population density, elevation) from a dynamic branch (multi-head self-attention
over geographic relationships and inter-pollutant coupling), then treats inference as
residual correction anchored to multiple stations, fused by distance-weighted averaging
across stations. That last step is the substantive contribution: it removes the
Voronoi-partition discontinuities that single-anchor schemes produce, yielding a continuous
field. Reported as superior on large-scale real-world data with zero-shot cross-domain
generalisation, but the abstract names no dataset, no station count, no RMSE and no
baseline, so none of it is checkable here. The architecture is directly applicable to
dense low-cost networks, where the anchor stations are the few reference-grade nodes.}
{10.1016/j.neunet.2026.109575}

\paperentry{Wildfire CHON PM2.5 2026}
{FT-ICR MS with rule-based prediction and QSAR for nitrogen-containing organics in
wildfire PM$_{2.5}$}
{Environ Res}
{Molecular-level characterisation of PM$_{2.5}$ from a temperate coniferous forest wildfire
in South Korea, wildfire versus post-wildfire periods. GC-MS gave nitrated polycyclic
aromatic compounds \textbf{3.4$\times$ higher during the fire (0.122 vs 0.036\,ng/m$^3$)},
with 2-nitrofluorene then 3-nitrofluoranthene dominant; FT-ICR MS showed a relative shift
toward CHON formulae, and rule-based SOA prediction from 12 VOC precursors matched
limonene-derived products most often. QSAR LC$_{50}$ values of 1.0, 1.6 and 8.3\,mg/L for
$\beta$-pinene-, $\alpha$-pinene- and isoprene-derived compounds are ranked hazard, not
measured toxicity, and the authors say so --- the assignments are formula-level, so isomers
with very different toxicity are pooled. One fire, one forest type; the compositional shift
is the result, the hazard ranking is a hypothesis.}
{10.1016/j.envres.2026.125576}

\band{Occupational \& indoor exposure\hfill 1 record}{Clay}

\paperentry{Tangshan CWP IGF1 2026}
{Dust exposure, IGF1 and autophagy in coal workers' pneumoconiosis}
{Front Public Health}
{Retrospective analysis of 2{,}028 employees at one Tangshan coal enterprise (2023--2024)
with a paired SiO$_2$-stimulated macrophage model. Logistic regression retained early
initial working age, dust-exposure duration, smoking, white cell count and AST/ALT ratio;
patients and stimulated macrophages both showed raised IGF1, altered LC3/P62 and increased
PI3K/AKT/mTOR phosphorylation. Two weaknesses matter. Exposure is job history, not a
measured respirable dust concentration, so no exposure--response can be estimated; and the
cell model uses crystalline silica while the cohort's hazard is coal mine dust --- these are
different particles with different mechanisms, and the paper's own language stays
correlative. The white-cell and liver-enzyme ``risk factors'' are as plausibly consequences
of disease as causes.}
{10.3389/fpubh.2026.1899729}

\band{Mechanistic toxicology\hfill 1 record}{Violet}

\paperentry{Retinal CRM 2026}
{Certified particulate reference materials in primary mouse retinal cells}
{Biochem Biophys Rep}
{Two certified environmental particulate reference materials applied to primary mouse
retinal cells, plus repeated ocular-surface exposure in mice. Effects on DCFDA fluorescence
and ATP-dependent viability were condition-dependent and concentrated in glial conditions;
neither $\alpha$-tocopherol nor N-acetylcysteine gave a significant rescue, and glial
coculture did not improve retinal ganglion cell survival. In postnatal mice topical
particulate was associated with inner retinal thinning; in adult Thy1-CFP mice the
within-animal comparison was null and only a \emph{post hoc} between-animal pooling reached
significance (\textbf{--14.4 cells, --28.1 to --0.7, $P=0.040$}). The authors state plainly
that this establishes neither particle transport to the retina nor a causal oxidative
mechanism nor any ambient risk --- an unusually disciplined set of negative caveats, and the
reason this is tier C rather than a finding. Ocular-surface dosing is not inhalation, and
certified reference material is not ambient PM.}
{10.1016/j.bbrep.2026.102758}

\band{Burden, policy \& mitigation\hfill 2 records}{Coral}

\paperentry{CLHLS exercise 2026}
{Exercise, county industrial air pollution and all-cause mortality in older adults}
{Front Public Health}
{16{,}954 CLHLS participants aged 65+, 23 provinces, three waves over six years, Cox models
with time-varying covariates and sequential adjustment. Habitual exercise
\textbf{HR 0.835 (0.773--0.902)}; county industrial SO$_2$ \textbf{HR 1.001 (0.874--1.146)}
after adjustment; soot indicators null; and \textbf{no significant interaction} between
exercise and any of the four pollution indicators. The null is about the exposure metric as
much as the hypothesis: annual county soot and SO$_2$ emission and removal tonnage is an
administrative aggregate, not a concentration, and it carries neither PM$_{2.5}$ nor any
within-county gradient. Exercise is self-reported and reverse causation from subclinical
illness is not addressed. Reported here as a bounded null, not as evidence against
interaction.}
{10.3389/fpubh.2026.1894835}

\paperentry{CA wine-grape smoke 2026 (preprint)}
{Quantifying impacts of wildfire smoke on crop yields in California}
{arXiv (preprint)}
{Statistical regression of wine-grape yield on particulate matter alongside climate
predictors at three Northern California sites. PM entered as a significant predictor
jointly with temperature at two of the three, with final models reaching
\textbf{$R^2=0.79$ at Sonoma and 0.66 at Napa} --- read by the authors as a shift from
climate-controlled to wildfire-and-climate-controlled yields. Three sites, yield as an
annual aggregate, and no separation of smoke PM from background or from the heat and
radiation anomalies that accompany fire weather; with that many candidate predictors, an
$R^2$ of this size is weak evidence of a specific PM pathway. Included as the only record
this window treating PM as an economic rather than a health exposure, and because the
smoke-attribution problem it dodges is the same one source-resolved monitoring exists to
solve. Not peer reviewed.}
{10.48550/arXiv.2609.04422}

% ==================================================================== TABLE
\section*{\sffamily\bfseries\color{Deep}Corpus register}
\vspace{-1.5mm}{\color{Amber}\rule{\linewidth}{1.1pt}}\vspace{2.5mm}

{\fontsize{6.7}{8.4}\selectfont
\setlength{\tabcolsep}{2.4pt}
\begin{longtable}{@{}>{\RaggedRight}p{25mm}>{\RaggedRight}p{41mm}>{\RaggedRight}p{27mm}%
>{\RaggedRight}p{23mm}>{\RaggedRight}p{18mm}>{\RaggedRight}p{15mm}@{}}
\toprule
\textsf{\textbf{First author}} & \textsf{\textbf{Focus}} & \textsf{\textbf{Journal}} &
\textsf{\textbf{Design}} & \textsf{\textbf{Metric}} & \textsf{\textbf{Region}}\\
\midrule
\endfirsthead
\toprule
\textsf{\textbf{First author}} & \textsf{\textbf{Focus}} & \textsf{\textbf{Journal}} &
\textsf{\textbf{Design}} & \textsf{\textbf{Metric}} & \textsf{\textbf{Region}}\\
\midrule
\endhead
\multicolumn{6}{@{}l@{}}{\cellcolor{Mist}\textsf{\textbf{\textcolor{Deep}{Neuro / mental health}}}}\\[0.6mm]
\href{https://doi.org/10.1001/jamanetworkopen.2026.25015}{Dales} & Wildfire-sourced PM$_{2.5}$ and ED visits for migraine and other primary headache & \textit{JAMA Netw Open} & Case-crossover & PM$_{2.5}$ (wildfire) & Canada\\
\href{https://doi.org/10.1093/ageing/afag215}{Glover} & Umbrella meta-review of urban environmental exposures and cognitive health & \textit{Age Ageing} & Umbrella review / meta-analysis & PM$_{2.5}$ & Global\\
\href{https://doi.org/10.1016/j.nsa.2026.107019}{Goossens} & In utero air pollution exposure and pre-/postnatal brain development & \textit{Neurosci Appl} & Narrative review & PM$_{2.5}$ / PM$_{10}$ / PM$_{\mathrm{coarse}}$ & Global\\
\href{https://doi.org/10.1016/j.isci.2026.116617}{Ma} & PM$_{2.5}$-induced neurotoxicity and depression: exposure risks and mechanisms & \textit{iScience} & Narrative review & PM$_{2.5}$ + components & Global\\
\href{https://doi.org/10.3760/cma.j.cn112150-20250804-00752}{Pan} & Maternal PM exposure and offspring neurodevelopment: effects and mechanisms & \textit{Zhonghua Yu Fang Yi Xue Za Zhi} & Narrative review & PM (bulk) & China\\
\addlinespace[0.9mm]
\multicolumn{6}{@{}l@{}}{\cellcolor{Mist}\textsf{\textbf{\textcolor{Teal}{Cardiovascular \& metabolic}}}}\\[0.6mm]
\href{https://doi.org/10.3760/cma.j.cn112338-20251120-00843}{Huang} & PM$_{2.5}$ and dyslipidaemia in rural Chinese children; dietary effect modification & \textit{Zhonghua Liu Xing Bing Xue Za Zhi} & Cross-sectional (multistage) & PM$_{2.5}$ & China\\
\href{https://doi.org/10.32394/pe/218966}{Kazmierska} & Air pollution and cardiovascular disease: review of current data & \textit{Przegl Epidemiol} & Narrative review & PM$_{2.5}$ / PM$_{10}$ & Global\\
\href{https://doi.org/10.1007/s00125-026-06814-2}{Xi} & Air pollution, built environment and temperature vs. prediabetes phenotypes & \textit{Diabetologia} & Prospective cohort & PM$_{2.5}$ & Germany\\
\href{https://doi.org/10.1161/JAHA.125.046917}{Yang} & Short-term PM$_{2.5}$ and hospitalisation/ED use in atrial fibrillation & \textit{J Am Heart Assoc} & Case-crossover & PM$_{2.5}$ (wildfire-influenced) & USA\\
\addlinespace[0.9mm]
\multicolumn{6}{@{}l@{}}{\cellcolor{Mist}\textsf{\textbf{\textcolor{Sky}{Reproductive \& developmental}}}}\\[0.6mm]
\href{https://doi.org/10.1016/j.jes.2025.10.043}{Chen} & Prenatal air pollution, fetal growth trajectory phenotypes and child neurodevelopment & \textit{J Environ Sci (China)} & Prospective cohort & PM$_{2.5}$ / PM$_{10}$ & China\\
\href{https://doi.org/10.1016/j.ecoenv.2026.120566}{Huang} & Joint prenatal bisphenol + air-pollutant mixture exposure and preschool IQ & \textit{Ecotoxicol Environ Saf} & Prospective cohort (mixtures) & PM$_{2.5}$ / PM$_{10}$ & China\\
\href{https://doi.org/10.1007/s40572-026-00552-8}{Peterson} & Ambient air pollution, PAHs and pubertal development: critical review & \textit{Curr Environ Health Rep} & Critical review & PM$_{2.5}$ components & Global\\
\href{https://doi.org/10.1016/j.envres.2026.125313}{Xie} & Long-term PM$_{2.5}$ and timing of pubertal development in boys and girls & \textit{Environ Res} & Prospective cohort & PM$_{2.5}$ & China\\
\addlinespace[0.9mm]
\multicolumn{6}{@{}l@{}}{\cellcolor{Mist}\textsf{\textbf{\textcolor{Sage}{Respiratory \& allergic}}}}\\[0.6mm]
\href{https://doi.org/10.1007/s00484-026-03289-z}{Bayar Muluk} & Climate change and otolaryngology: airway disease impacts and emission reduction & \textit{Int J Biometeorol} & Narrative review & PM$_{2.5}$ & Global\\
\href{https://doi.org/10.4103/mgr.MEDGASRES-D-25-00164}{Liao} & Short-term air pollution, respiratory symptoms and lung function in asthma & \textit{Med Gas Res} & Panel study & PM$_{2.5}$ / PM$_{10}$ & China\\
\href{https://doi.org/10.3390/epidemiologia7040097}{Sanchez de Prada} & Ambient PM and asthma-coded admissions in children 0-3 y, Spain & \textit{Epidemiologia} & Ecological panel & PM$_{2.5}$ / PM$_{10}$ & Spain\\
\href{https://doi.org/10.1016/j.jhazmat.2026.143076}{Yang} & Allergenic vs. toxic risk of airborne seed hairs compared with ambient PM$_{2.5}$ & \textit{J Hazard Mater} & Field measurement & PM$_{2.5}$ + bioaerosol & China\\
\addlinespace[0.9mm]
\multicolumn{6}{@{}l@{}}{\cellcolor{Mist}\textsf{\textbf{\textcolor{Violet}{Mechanistic toxicology}}}}\\[0.6mm]
\href{https://doi.org/10.1016/j.jes.2025.12.074}{Ding} & PM$_{2.5}$ promotes vascular calcification via miR-27b-5p/HES1 signalling & \textit{J Environ Sci (China)} & In vitro + organoid + murine & PM$_{2.5}$ & China\\
\href{https://doi.org/10.1007/s00360-026-01698-5}{Maraschi} & Osmoregulatory disruption and metal bioaccumulation from settleable atmospheric PM & \textit{J Comp Physiol B} & Controlled exposure (ecotox) & Settleable APM & Brazil\\
\href{https://doi.org/10.1016/j.jes.2025.11.018}{Wu} & Macrophage NLRP3 inflammasome in real-ambient PM-induced hepatic glucose dysregulation & \textit{J Environ Sci (China)} & Real-ambient murine + co-culture & Real-ambient PM & China\\
\addlinespace[0.9mm]
\multicolumn{6}{@{}l@{}}{\cellcolor{Mist}\textsf{\textbf{\textcolor{Clay}{Occupational \& indoor}}}}\\[0.6mm]
\href{https://doi.org/10.3389/fmedt.2026.1796828}{Chen} & Nested suction instrument to cut intraoperative surgical-smoke PM$_{2.5}$ & \textit{Front Med Technol} & Bench / device experiment & PM$_{2.5}$ (surgical smoke) & China\\
\href{https://doi.org/10.1093/annweh/wxag055}{Cole} & Determinants of respirable crystalline silica exposure in Australian tunnelling & \textit{Ann Work Expo Health} & Observational exposure study & Respirable dust / RCS & Australia\\
\href{https://doi.org/10.1016/j.jhazmat.2026.143081}{Ding} & Stage-resolved hair Zn-isotope dynamics as a firefighting combustion-exposure tracer & \textit{J Hazard Mater} & Repeated-measures biomarker & Combustion-derived PM & China\\
\href{https://doi.org/10.3390/jox16040119}{Ramsden} & Xenobiotic hazards in aircraft cabin air & \textit{J Xenobiot} & Review + risk economics & Ultrafine particles & Global\\
\addlinespace[0.9mm]
\multicolumn{6}{@{}l@{}}{\cellcolor{Mist}\textsf{\textbf{\textcolor{Amber}{Sensing, forecasting \& instrumentation}}}}\\[0.6mm]
\href{https://doi.org/10.1016/j.jenvman.2026.130515}{Shi} & Adaptive spatiotemporal graph transformer for multi-site multi-step PM$_{2.5}$ forecasting & \textit{J Environ Manage} & Deep-learning model & PM$_{2.5}$ & China / India\\
\href{https://doi.org/10.3390/bios16070380}{Tsuchiya} & Seasonal and indoor-outdoor characteristics of fluorescent aerosol particles in offices & \textit{Biosensors} & Real-time sensor campaign & Bioaerosol / coarse & Japan\\
\addlinespace[0.9mm]
\multicolumn{6}{@{}l@{}}{\cellcolor{Mist}\textsf{\textbf{\textcolor{Amber}{Exposure assessment \& modelling}}}}\\[0.6mm]
\href{https://doi.org/10.1029/2025GH001621}{Solorzano-Araque} & Multiscale exposure modelling: LOTOS-EUROS + mobility surveys + particle toxicity & \textit{GeoHealth} & CTM + mobility integration & PM$_{2.5}$ / PM$_{10}$ & Colombia\\
\addlinespace[0.9mm]
\multicolumn{6}{@{}l@{}}{\cellcolor{Mist}\textsf{\textbf{\textcolor{Coral}{Burden, policy \& mitigation}}}}\\[0.6mm]
\href{https://doi.org/10.1002/hsr2.72886}{Karamzad} & Ischaemic heart disease burden in MENA 1990-2021; PM attributable fraction 28.2\% & \textit{Health Sci Rep} & GBD secondary analysis & Ambient PM & MENA\\
\href{https://doi.org/10.1371/journal.pone.0347989}{Rees} & HEPA filtration to reduce respiratory infection in care homes & \textit{PLoS One} & RCT process evaluation & Filtered PM & UK\\
\href{https://doi.org/10.1016/j.jes.2026.01.009}{Shan} & Policy-driven emission reductions and post-pandemic rebound, high-resolution inventories & \textit{J Environ Sci (China)} & Emission inventory & PM emissions & China\\
\addlinespace[0.9mm]
\multicolumn{6}{@{}l@{}}{\cellcolor{Mist}\textsf{\textbf{\textcolor{Slate}{Other clinical endpoints}}}}\\[0.6mm]
\href{https://doi.org/10.3390/dj14070440}{Lupita} & PM$_{2.5}$, preventive-care capacity and caries/edentulism burden across the EU & \textit{Dent J} & Ecological panel & PM$_{2.5}$ & EU-27\\
\href{https://doi.org/10.1016/j.jdin.2026.06.008}{Park} & Long-term PM exposure and risk of androgenetic alopecia & \textit{JAAD Int} & Nationwide cohort & PM$_{2.5}$ / PM$_{10}$ & South Korea\\
\href{https://doi.org/10.3171/2026.3.PEDS2615}{Sun} & Environmental exposure and area deprivation vs. paediatric hydrocephalus outcomes & \textit{J Neurosurg Pediatr} & Registry cohort + geospatial & PM$_{2.5}$ (satellite-derived) & USA\\
\addlinespace[0.9mm]
\bottomrule
\end{longtable}}


\vspace{4mm}
\begin{keybox}
{\sffamily\bfseries\footnotesize\color{Deep}Provenance and method}\\[1.2mm]
{\fontsize{7.2}{8.8}\selectfont
Entry window \textbf{6 September 2026}, a single day, run against \texttt{[Date - Entry]}.
\textbf{PubMed} (NCBI E-utilities) returned \textbf{9 health-axis} and \textbf{1
sensing-axis} hits to the harvester; the \textbf{PubMed connector}, run independently on
both axes, returned \textbf{10} and \textbf{13}, contributing three in-scope records the
harvester missed. The sensing-axis surplus is almost entirely false positive --- free-text
\texttt{"remote sensing"} and \texttt{"exposure assessment"} matched marine phycotoxins,
sports biomechanics, volcanic petrology and MODIS evapotranspiration --- and \textbf{10 of
those 13 were rejected}, a third in-scope hit being one the health axis had already
returned. \textbf{Europe PMC} (\texttt{CREATION\_DATE}) returned 6 raw,
\textbf{Crossref-by-ISSN} 11 raw, \textbf{arXiv} 1 inside the local recency screen.
\textbf{OpenAlex} returned HTTP 429, unchanged since 28 July. \textbf{Consensus} returned
three hits, all out of window on Crossref \texttt{created} date (2019, 2022, 2023) --- a nil
for the eighth consecutive run, and the absence of a date post-filter remains the cause.
\textbf{Scholar Gateway} was queried but its response exceeded the reply budget and was not
consumed; the abstracts already carried the quantities needed, so no record depended on it.
\textbf{ClinicalTrials.gov} posted \textbf{no updates registry-wide} on 6 September ---
verified with a no-term control, as on 5 September and 29/30 August --- so the trial zero is
a posting-calendar artefact of the weekend and is not evidence about PM trial activity.
Twenty-eight raw records deduplicated to 23 unique on PMID $\to$ DOI $\to$ normalised title;
none had been shipped in a prior issue. Twenty-one records were screened out with a logged
reason, chiefly non-atmospheric \textit{Environ Sci Technol} chemistry, built-environment
and energy papers with no particle metric, ozone- and VOC-only studies, one corrigendum,
and the connector's sensing-axis false positives.
Subtopic, design, particle-metric and geography labels are assigned from title, abstract
and keywords; assignment is single-label by dominant contribution and is a judgement, not a
controlled vocabulary. \textit{Global / multi-region} in the geography panel is a fallback
bucket that absorbs unstated and non-geographic settings. Effect estimates are transcribed
verbatim and use heterogeneous contrasts --- see the forest-plot caption.
All DOI links resolve to the publisher of record; the arXiv record is cited by its DataCite
DOI. Abstracts remain the copyright of their respective publishers.\par}
\end{keybox}

\end{document}
